\answer{ex:19:1}{(a)~Una sinapsis da $6.7\mV$ ($R=66.7$~M$\Omega$), así que para llegar al umbral en absoluto hacen falta al menos $4$ (tres bastan solo asintóticamente); $8$ en $10\ms$; $14$ en $5\ms$. (b)~$0.1\ms\ll\tau_m$; la corrección por la fuga es de $\sim0.25\%$.}
\answer{ex:19:2}{(a)~$1.75\cdot10^{10}$ (una cuarta parte de los mezcladores responde a todo), $4.4\cdot10^9$ y $0.06$ (con $k=40$ casi nunca responde ninguno). (b)~$10^3$ veces: $2\cdot10^5$ frente a $200$.}
\answer{ex:19:3}{(a)~$C(2n,n)=2^{2n-1}$ por la simetría de los coeficientes binomiales. (b)~$0.93$ y $0.09$; la anchura de la transición es $\sim\sqrt{2n}$ en $P$, y la relativa, $\sim1/\sqrt n$.}
\answer{ex:19:4}{Desde una sola dendrita, $V_{\text{cuerpo}}<\kappa E_E<\theta$ con cualquier número de entradas; con $g_0=0.5g_L$ hacen falta $n\ge3$ en cada dendrita.}
\answer{ex:19:5}{$I\ge C\sigma_V/\sigma_t=15$~nA, es decir, $150$ sinapsis síncronas: poco realista; a la neurona pequeña le bastaría $1$~nA, y con una sinapsis de $4$~nA la fluctuación es de $5$~\textmu s.}
\answer{ex:19:6}{Máximo en el extremo cercano: $0.03\,\tau_m$ y $0.97$ ($L=0.2$); $0.32\,\tau_m$ y $0.66$ ($L=1$); $0.79\,\tau_m$ y $0.33$ ($L=2$). La estimación~\eqref{eq:dendriteload} con la capacidad total solo vale para $L\ll1$.}
\answer{ex:19:7}{Ejercicio abierto. Con $m$ fijo, el tiempo de respuesta crece linealmente con el $P$ memorizado; con una fracción fija $m=\varphi n$ deja de depender de $n$, y la lentitud de la neurona grande es consecuencia de sumar muchas entradas, no del tamaño en sí.}
